\documentclass[9pt,letterpaper,twocolumn]{article}

\usepackage[T1]{fontenc}
\usepackage{lmodern}
\usepackage[margin=0.62in,columnsep=0.27in]{geometry}
\usepackage{amsmath,amssymb,mathtools,bm}
\usepackage{array,booktabs}
\usepackage{xcolor}
\usepackage{microtype}
\usepackage{fancyhdr}
\usepackage{needspace}

% The user-local tools/array package is newer than the available LaTeX kernel.
\makeatletter
\def\vcenter@text{\vbox}
\makeatother

\definecolor{navy}{HTML}{17365D}
\definecolor{green}{HTML}{2F6B4F}

\pagestyle{fancy}
\fancyhf{}
\fancyhead[L]{\footnotesize Holstein-dimer stability addendum}
\fancyhead[R]{\footnotesize New print supplement}
\fancyfoot[C]{\thepage}
\renewcommand{\headrulewidth}{0.3pt}

\setlength{\parindent}{0pt}
\setlength{\parskip}{2.2pt plus 0.4pt}
\renewcommand{\arraystretch}{1.07}

\newcommand{\R}{\mathbb R}
\newcommand{\I}{\mathbb I}
\newcommand{\ii}{\mathrm i}
\newcommand{\segmenthead}[1]{%
  \Needspace{4\baselineskip}%
  \par\smallskip{\color{navy}\large\bfseries #1}\par
  \vspace{-2pt}{\color{navy}\rule{\columnwidth}{0.7pt}}\par\smallskip
}
\newcommand{\conceptbreak}{%
  \par\smallskip\noindent
  {\color{black!30}\rule{\columnwidth}{0.35pt}}%
  \par\smallskip
}

\begin{document}
\raggedbottom
\setcounter{page}{24}

\twocolumn[
\begin{center}
  {\color{navy}\Large\bfseries New Print Supplement}\par\vspace{2pt}
  {\large Minimum-Norm Solver and Offline Closure Diagnostic}\par\vspace{2pt}
  {\small Continuation after the previously printed page 23, 2 August 2026}
\end{center}
\vspace{2pt}\hrule\vspace{5pt}
]

\segmenthead{Supplement to E. How Eq. (X23) was solved}

At one right-hand-side evaluation, the printed correction problem chooses
\(w\in\R^{31}\).  Its three linear equalities are energy neutrality and the
real and imaginary parts of the fixed-number correlation-trace condition.
Collect them as
{\footnotesize
\[
Lw=d,
\qquad
L=
\begin{pmatrix}
\nabla_xE_{\rm int}(x)^{\mathsf T}\\
e_{17}^{\mathsf T}\\
e_{18}^{\mathsf T}
\end{pmatrix},
\qquad
d=
\begin{pmatrix}
0\\
-\dot x_{0,17}-\beta x_{17}\\
-\dot x_{0,18}-\beta x_{18}
\end{pmatrix}.
\tag{X26e}
\]
}
Here \(e_j\) is the \(j\)-th standard coordinate vector.  A singular-value
decomposition gives the minimum-Euclidean-norm particular solution
\(w_{\rm p}=L^+d\) and an orthonormal basis
\(Z\in\R^{31\times(31-r)}\) for the null space of \(L\):
\[
Lw_{\rm p}=d,
\quad LZ=0,
\quad Z^{\mathsf T}Z=I,
\quad Z^{\mathsf T}w_{\rm p}=0.
\tag{X26f}
\]
Every equality-feasible correction is therefore
\[
\begin{aligned}
w&=w_{\rm p}+Zy,
&y&\in\R^{31-r},\\
\frac12\lVert w\rVert_2^2
&=\frac12\lVert w_{\rm p}\rVert_2^2
+\frac12\lVert y\rVert_2^2.
\end{aligned}
\tag{X26g}
\]
The equalities fix \(w_{\rm p}\).  The remaining minimization is the Euclidean
projection of \(y=0\) onto the semidefinite-feasible set.

\conceptbreak
\textbf{One negative eigenvector produces one linear cut.}
Let \(\ell\in\{\rho^-,\rho^+,\mathcal G\}\) select the electronic lower
barrier, electronic upper barrier, or joint Gram barrier.  After equality
elimination, its corrected barrier matrix is the affine map
\[
B_\ell(y)=M_{0,\ell}+\Delta M_\ell(w_{\rm p}+Zy).
\tag{X26h}
\]
At the current reduced point \(y^{(k)}\), diagonalization exposes a normalized
failing direction \(v\):
\[
B_\ell(y^{(k)})v=\lambda v,
\qquad v^\dagger v=1,
\qquad \lambda<0.
\tag{X26i}
\]
The eigenvector \(v\) identifies the failing matrix direction, and
\(\lambda\) measures its differential-barrier violation.  The response of
that direction to the \(j\)-th correction coordinate is
\[
[a_\ell(v)]_j=v^\dagger\Delta M_\ell(e_j)v,
\qquad n=Z^{\mathsf T}a_\ell(v).
\tag{X26j}
\]
Because the lift is linear, the exposed quadratic form is exactly affine:
\[
\begin{aligned}
v^\dagger B_\ell(y)v
&=\lambda+n^{\mathsf T}(y-y^{(k)}),\\
v^\dagger B_\ell(y)v\geq0
&\Longleftrightarrow
n^{\mathsf T}y\geq n^{\mathsf T}y^{(k)}-\lambda.
\end{aligned}
\tag{X26k}
\]
Thus one negative eigenvector supplies one half-space containing every
correction that is feasible in that matrix direction.  The vector \(n\) is
normal to its boundary hyperplane.

For a single cut \(a^{\mathsf T}w+\lambda\geq0\), the minimum-norm solution
follows directly:
{\footnotesize
\[
\begin{aligned}
&\underset{w}{\operatorname{minimize}}\ \tfrac12w^{\mathsf T}w
\quad\text{subject to}\quad a^{\mathsf T}w+\lambda\geq0,\\
&\mathcal L(w,\mu)=\tfrac12w^{\mathsf T}w
-\mu(a^{\mathsf T}w+\lambda),\\
&\nabla_w\mathcal L=w-\mu a=0,
\qquad
\mu=-\lambda/\lVert a\rVert_2^2,\\
&\boxed{w^\star=-\lambda a/\lVert a\rVert_2^2.}
\end{aligned}
\tag{X26d}
\]
}
This is the Euclidean projection of the origin onto the exposed boundary.

\conceptbreak
\textbf{Accumulated cuts give the implemented quadratic program.}
After \(p\) failing directions have been exposed, the reduced problem is
\[
\boxed{
\begin{aligned}
\underset{y\in\R^{31-r}}{\operatorname{minimize}}\quad
&\tfrac12y^{\mathsf T}y\\
\text{subject to}\quad
&n_i^{\mathsf T}y\geq c_i,
\qquad i=1,\ldots,p.
\end{aligned}}
\tag{X26l}
\]
SciPy's Sequential Least Squares Programming optimizer, SLSQP, projects the
reduced origin onto the intersection of these accumulated half-spaces.  The
full correction and matrices are then reconstructed and diagonalized again:
\[
y^{(k)}
\longrightarrow B_\ell(y^{(k)})
\longrightarrow(\lambda,v)
\longrightarrow(n^{\mathsf T}y\geq c)
\longrightarrow y^{(k+1)}.
\tag{X26m}
\]
Previously generated cuts remain active.  Iteration stops when all three
complete barrier matrices satisfy the cone tolerance.  Geometrically, the
positive-semidefinite feasible set is the intersection of the half-spaces
generated by every normalized matrix direction; the accumulated finite set is
a cutting-plane approximation to that spectrahedron.

Near an eigenvalue degeneracy, the exposed lowest eigenvector can exchange
with another direction.  The fallback asks SLSQP to impose every eigenvalue
constraint directly:
\[
\begin{aligned}
\underset{y\in\R^{31-r}}{\operatorname{minimize}}\quad
&\tfrac12y^{\mathsf T}y\\
\text{subject to}\quad
&\lambda_m(B_\ell(y))\geq0
\quad\text{for every }\ell,m,
\end{aligned}
\tag{X26n}
\]
with the simple-eigenvalue derivative
\[
\frac{\partial\lambda_m(B_\ell(y))}{\partial y_j}
=v_m^\dagger\frac{\partial B_\ell(y)}{\partial y_j}v_m.
\tag{X26o}
\]
At the exact cutoff-16 initial state, the direct and cutting-plane solves
agreed in correction norm to relative tolerance \(2\times10^{-7}\).  Both
minimize the Euclidean norm of the independent 31-coordinate correction.
Minimizing the Frobenius norm of all displayed lifted matrix entries requires
the weighted metric derived later in Eq. (X48).

\segmenthead{Offline reconstruction of the missing correlation velocity}

The pair
\[
C=(C^0,C^1),
\qquad C^q\in\mathbb C^{2\times2},
\]
contains one electron--phonon correlation matrix for each phonon mode.  The
symbol \(\dot C_{31}\) denotes the derivative assigned to this pair by the
31-real-coordinate archive closure:
\[
x_{\rm ex}(t)
\xrightarrow{\mathcal E}X_{\rm ex}(t)
\xrightarrow{F_C}\dot C_{31}(t).
\tag{X44a}
\]
The closure is evaluated on the exact contracted state at each saved time;
the state used in this audit comes from the independently propagated truncated
Holstein Hamiltonian.  Printed Eq. (D14d) gives
\[
\dot C_{31}^{,q}
=-\ii\bigl(h(t)C^q-C^qh(t)+\omega_qC^q\bigr)
+\sum_{a=0}^{4}S_a^q.
\tag{X44b}
\]

The exact operator derivative contains three additional physical sources.
With \(O_{ij}=c_{j\uparrow}^\dagger c_{i\uparrow}\),
\(n_{r\uparrow}=c_{r\uparrow}^\dagger c_{r\uparrow}\), and
\(\delta X_r=\delta b_r+\delta b_r^\dagger\), define
{\footnotesize
\[
\begin{aligned}
Q^{qr}_{ij}
&=\left\langle[O_{ij},n_{r\uparrow}]\delta X_r\delta b_q\right\rangle,\\
Q^{qr}_{ij,0}
&=\left\langle[O_{ij},n_{r\uparrow}]\right\rangle
  \left\langle\delta X_r\delta b_q\right\rangle,\\
K^{qr}_{ij}&=Q^{qr}_{ij}-Q^{qr}_{ij,0},\\
P^q_{ij}&=\delta_{iq}\rho_{qj}-\rho_{ij}\rho_{qq},\\
D^q_{ij}&=\left\langle O_{ij}n_{q\downarrow}\right\rangle
-\rho_{ij}\left\langle n_{q\downarrow}\right\rangle.
\end{aligned}
\tag{X44c}
\]
}
Their velocity additions are
{\footnotesize
\[
\begin{aligned}
\Delta\dot C^q_{K,ij}
&=-\ii g\sum_rK^{qr}_{ij},\\
\Delta\dot C^q_{P,ij}
&=-\ii gP^q_{ij}
-\left[\left(\sum_{a=0}^{4}S_a^q\right)_{ij}
+\ii g\sum_rQ^{qr}_{ij,0}\right],\\
\Delta\dot C^q_{D,ij}
&=-\ii gD^q_{ij}.
\end{aligned}
\tag{X44d}
\]
}
The \(P\) line is a replacement: it subtracts the effective same-spin source
already contained in the archive factorization and supplies the exact
fixed-number Pauli covariance.

At every exact-trajectory sample, the audit formed
\[
\begin{aligned}
&\dot C_{31},\\
&\dot C_{31}+\Delta\dot C_K,\\
&\dot C_{31}+\Delta\dot C_K+\Delta\dot C_P,\\
&\dot C_{31}+\Delta\dot C_K+\Delta\dot C_P+\Delta\dot C_D.
\end{aligned}
\]
and compared each expression with \(\dot C_{\rm ex}\).  Over
\(0\leq t\leq4\), the residual-subtracted 14-real-coordinate errors were
\[
\begin{array}{@{}lrr@{}}
\toprule
\text{supplied sources}&\operatorname{RMS}\lVert e\rVert_2
&\max_t\lVert e\rVert_2\\
\midrule
\varnothing&0.166690&0.274310\\
K&0.089345&0.137864\\
K+P&0.026160&0.038268\\
K+P+D&2.53\times10^{-5}&3.47\times10^{-5}\\
\bottomrule
\end{array}
\tag{X44e}
\]
The remaining scale matches the finite-cutoff commutator remainder
\[
R_{\rm cut}
=\dot C_{\rm ex}-\dot C_{31}
-\Delta\dot C_K-\Delta\dot C_P-\Delta\dot C_D.
\tag{X44f}
\]

This sequence was an offline local-derivative diagnostic.  The exact
trajectory supplied \(K\), \(P\), and \(D\) only for attribution.  These
terms were never supplied to RK4 or to the minimum-norm controller \(w\).
The controller preserves representability; the offline audit identifies the
missing dynamics that an enlarged autonomous closure must reproduce.

\end{document}
